\subsection{A trivial middle action and a binary adjacent stalk}

\begin{lemma}[Proved: collapsing a middle corner with trivial outgoing action]
\label{CP2:A-linear-trivial-middle-action-H2}
Let $C$ have full linear support $0<1<2$, finite corners
$A,B,N$, and nonempty forward hom-sets $X=C(0,1)$,
$Y=C(1,2)$ and $Z=C(0,2)$.  Suppose the right $B$ action
on $Y$ is trivial.  There is an actual quotient functor
\[
 F:C\longrightarrow \overline C,
 \qquad
 (A,B,N;X,Y,Z)\longmapsto(A,1,N;B\backslash X,Y,Z).
\]
Its pullback is an isomorphism on rational $H^2$.
No projectivity or homological boundedness hypothesis is required.
\end{lemma}
\begin{proof}
Composition descends because
$\mu(y,bx)=\mu(yb,x)=\mu(y,x)$; the right $A$ action
on $X$ descends to its $B$ orbit set as well.  All other
products are unchanged, so this is a literal category functor.

We give the low-degree comparison explicitly.  All following
complexes are normalized rational bar complexes in nonnegative
homological degrees.  Set
\[
 V=C_*(BB),\qquad
 U=C_*B(B,X,A),\qquad
 U_0=C_*B(B\backslash X,A),\qquad
 T=C_*B(N,Y).
\]
Here $U$ is the two-sided edge bar, with its $B$ and $A$
directions, and $U_0$ is the right $A$ action bar of the
orbit set.  Write
\[
 V^+=\ker(V\longrightarrow\mathbb Q),\qquad
 U^+=\ker(U\longrightarrow U_0).
\]
The two displayed maps are degreewise surjective.  Before
taking the $A$ bar, the map from the left $B$ action bar
of $X$ to $\mathbb Q(B\backslash X)$ has kernel with zero
zeroth homology.  The first-quadrant $A$-bar filtration
therefore gives $H_0(U^+)=0$.  A finite monoid has zero
rational first homology, so
\[
 H_0(V^+)=H_1(V^+)=0.
\]
The augmentation $T\to\mathbb Q$ is surjective on
zeroth homology because $Y$ is nonempty.

Use the strict-object bar model for a full linear category.
It consists of the corner complexes, the three edge complexes
shifted once, and the triangle complex shifted twice, with
their actual product-face maps.  The trivial right $B$ action
identifies the middle-to-last edge complex with $T\otimes V$
and the triangle complex with $T\otimes U$, by natural
Eilenberg--Zilber comparison.  These identifications respect
all product and augmentation faces.  In the relative kernel
of the comparison with $\overline C$, the only four terms are
\[
 V^+,\qquad U^+[1],\qquad
 (T\otimes V^+)[1],\qquad (T\otimes U^+)[2].
\]
The product face to the long-edge bar vanishes on $U^+$:
that product factors through $B\backslash X$.  The faces to
the unchanged first and last corner complexes vanish for
the same augmentation reason.  The remaining faces are
$U^+\to V^+$ and the augmentation of $T$.  With the usual
total-complex signs, this relative kernel is consequently
\[
 K\simeq
 \operatorname{Cone}(U^+\longrightarrow V^+)
 \otimes
 \operatorname{Cone}(T\longrightarrow\mathbb Q).
\]
This description may equivalently be made before the natural
Eilenberg--Zilber comparison, using its two bar directions;
it computes the homotopy fiber of the actual chain map $F$.
The first factor has no homology below degree two, and the
second has none below degree one.  Over $\mathbb Q$, Kunneth
therefore gives $H_j(K)=0$ for $j\leq2$.  The relative long
exact sequence makes $F_*:H_2(BC)\to H_2(B\overline C)$ an
isomorphism.  Rational duality proves the asserted statement,
and the comparison preserves products because $F$ is actual.
\end{proof}

\begin{theorem}[Proved: a singleton long hom-set and a binary adjacent hom-set]
\label{CP2:A-linear-singleton-long-binary-any-corners}
Let $C$ have full linear support, a singleton long hom-set,
and at least one adjacent hom-set of size at most two.
There is no restriction on its finite corners or its other
adjacent hom-set.  Every rational degree-two product vanishes.
\end{theorem}
\begin{proof}
By opposition suppose $|Y|\leq2$ in the notation above,
and write $Z=\{z\}$.  Put $D=C|_{\{0,1\}}$.

If the right $B$ action on $Y$ has a fixed state $y$, the
components $z,y$ define a natural transformation from the
inclusion of $D$ to the constant functor at object $2$:
$za=z$, $yb=y$, and $yx=z$ for every $a,b,x$.  The unique
long arrow also defines a natural transformation from the
constant functor at object $0$ to the inclusion of $N$.
Thus both vertex-block restrictions of every positive class
are zero.  The two-block relative suspension argument kills
their products.  This includes $|Y|=1$.

Assume $|Y|=2$.  If the right $B$ image contains both
constant transformations, a commuting left $N$ actor fixes
both states.  Thus $N$ acts trivially on $Y$ and on $Z$.
Lemma~\ref{CP2:A-trivial-cut-H2} gives an actual functor
replacing $N$ by $1$ whose pullback is an $H^2$ isomorphism.
The quotient is an act cone on $D$ with stalk sizes one and
two.  Lemma~\ref{CP2:A-high-affine} kills its positive
products, hence also all second products in $BC$.

If exactly one constant occurs in the right $B$ image,
the interchange cannot occur, since it would produce the
other constant.  Its value is a common fixed state, already
handled.  The remaining transformation images are the
identity image and the permutation group $C_2$.

For the identity image, the preceding lemma identifies
$H^2(BC)$ with the pullback from a category having trivial
middle corner and singleton long hom-set.  Its second
products vanish by
Lemma~\ref{CP2:A-trivial-middle-singleton-long}.

Finally suppose the interchange occurs.  Every commuting
left $N$ map of $Y$ is then a permutation.  The vector
$q=([y_0]+[y_1])/2$ is fixed by both corner actions.
Let $P$ be the connecting $(N,D)$ profunctor, whose stalks
at $0,1$ are $Z,Y$.  The assignments
\[
 1\longmapsto[z]\quad\text{at }0,
 \qquad 1\longmapsto q\quad\text{at }1
\]
give an $(N,D)$-linear section of the augmentation
$\mathbb QP\to\mathbf1$: every $x\in X$ sends both
members of $Y$ to $z$, so it sends $q$ to $[z]$.
Applying the two bar directions gives a chain section of
the correspondence projection to $BN\times BD$.
Consequently the joint map
\[
 H^2(BD;\mathbb Q)\oplus H^2(BN;\mathbb Q)
 \longrightarrow H^2(B(N,P,D);\mathbb Q)
\]
is injective.  Mayer--Vietoris forces both restrictions
of every collage class to vanish.  The same relative
suspension argument finishes this case.
\end{proof}

\begin{corollary}[Proved: three large linear nineteen-arrow profiles]
\label{CP2:A-linear19-three-binary-long-large-profiles}
The full-linear cardinality profiles
\[
 (1,4,8;3,2,1),\qquad
 (2,3,8;3,2,1),\qquad
 (3,2,8;3,2,1)
\]
cannot have classifying space of projective-plane type.
These are original indices $132,308,492$ in the frozen
nineteen-arrow linear inventory.  This assertion is
independent of any monoid enumeration.
\end{corollary}
\begin{proof}
Each has singleton long hom-set and a binary adjacent
hom-set.  Apply the theorem: a projective-plane generator
has nonzero rational square.
\end{proof}
